% Appendix, from exam/paper_b_solutions.md: the model answers to Paper B, with marks.

% The answers' headings are labelled ansb:q:N.
\renewcommand{\answerprefix}{ansb}

\chapter{Model Answers to Paper B}
\label{ansb:ch:answers}

Marks are shown per part. Method carries them: a correct derivation with a slip in it is worth more
than a correct answer with no working, and later parts consume earlier ones, so an error should be
followed through rather than penalised twice.

Where a part asks for VHDL, mark the \textbf{hardware described}, not the syntax. A listing with a
missing semicolon and the right circuit in it scores full marks; a listing that compiles and infers
a latch does not. The questions are in Appendix~\ref{exb:ch:paper}.

% ==========================================================================================
% The title is too long for one line: break it ragged, and before its last phrase.
\answersection{1}{A requirement, and the shape of the network it asks for}{12}

\answerpart{a}{2}
\begin{textdrawing}
run = left · right · curtain_clear · estop_n
\end{textdrawing}
\noindent\worth{1}

\textbf{Gates.} One four-input AND, or three two-input ANDs if the library has nothing wider.
\textbf{No inverter}, which is worth noticing before part (c) asks about it: the condition wanted is
``the stop has \emph{not} been pulled'', and that is \code{estop_n = 1} as it stands.

\textbf{The table.} Four inputs, sixteen rows, and exactly \textbf{one} of them is \code{1}: the row
where all four are \code{1}. \worth{1}

\answerpart{b}{4}
\begin{textdrawing}
run = estop_n · ( maint'·left·right·curtain_clear  +  maint·pedal )
\end{textdrawing}
\noindent\worth{1}

\textbf{The factored input: \code{estop_n}.} It multiplies the whole sum rather than appearing
inside either branch.

\textbf{What that expresses.} No mode can bypass the emergency stop. It is not an input to the mode
logic; it gates the mode logic's result. A third mode added next year cannot accidentally create a
path that runs the press with the stop pulled, because every path goes through the same AND.
\worth{1}

\textbf{The contrast.}

\begin{booktable}{@{}>{\raggedright\arraybackslash}p{8.4em}LL@{}}
  & \thead{Braking assistant} & \thead{Press control} \\
  \midrule
  Shape & \code{engine_brake = driver_brake + adas_brake} & \code{run = estop_n · (mode logic)} \\
  The override & is ORed in \textbf{last}, outside the fault logic & is ANDed in \textbf{last},
    outside the mode logic \\
  Forces the output & \textbf{high} & \textbf{low} \\
  Guarantees & the action is always \textbf{available}: the driver can always brake, including
    when the assistance system is faulty and doing nothing & the action can always be
    \textbf{prevented}: the press never cycles with the stop pulled, in any mode \\
\end{booktable}

\textbf{What routing the override through the logic would cost, in each case.} The same structural
mistake with opposite consequences, which is why the two are worth seeing together. Feed
\code{driver_brake} through the assistance system's fault logic and a fault stops the driver
braking --- the safety requirement the structure exists to protect. Feed \code{estop_n} into only
one branch of the sum, say the maintenance branch, and pulling the stop no longer stops the press
in normal mode. Both are one wire in the wrong place, and neither shows up in a test that only
exercises the happy path. \worth{1}

\textbf{The count.} Sixty-four rows, of which \textbf{10} are \code{1}:
\begin{itemize}
  \item \code{estop_n = 1} throughout, \code{maint = 0}: needs
    \code{left = right = curtain_clear = 1}, \code{pedal} free --- 2 rows.
  \item \code{estop_n = 1}, \code{maint = 1}: needs \code{pedal = 1}, the three others free --- 8
    rows.
\end{itemize}

\textbf{What the table would have cost you.} Ten AND-terms of six literals each, from a requirement
that is two clauses long, plus the whole job of simplifying them back down. The structure came out
of the \emph{statement} in one line. That is the case \secref{c1:a6} makes with the braking
assistant: derive from the requirement, and use the table to check the answer rather than to find
it. \worth{1}

\answerpart{c}{3}
\begin{vhdlcode}
library ieee;
use ieee.std_logic_1164.all;

entity press_control is
    port(left, right, curtain_clear, estop_n, maint, pedal: in  std_logic;
         run                                              : out std_logic);
end entity;

architecture behaviour of press_control is
-- The two mode branches, before the emergency stop gates either of them.
signal two_hand, maintenance: std_logic;
begin
    two_hand    <= (not maint) and left and right and curtain_clear;
    maintenance <= maint and pedal;

    run <= estop_n and (two_hand or maintenance);
end architecture;
\end{vhdlcode}

{\itshape (2 marks: 1 for the entity with the ports in order, 1 for an architecture matching the
candidate's own (b). The internal signals are optional; a single collapsed expression is equally
correct.)\par}

\medskip
\textbf{Where the \code{not} is, and where it is not.} The one inversion is on \code{maint}, an
\textbf{active-high} signal that has to be \emph{false} for the two-hand branch. \code{estop_n}, the
\textbf{active-low} signal, needs no inversion at all.

\textbf{The rule.} An \code{_n} suffix says what a \code{0} on that wire means; it does not say that
an inversion is needed. Write down the condition you want first, then invert whatever does not
already say it. ``Active-low, therefore \code{not}'' is the reflex to lose: here it would have
produced a press that runs only while the emergency stop is pulled. \worth{1}

\answerpart{d}{3}
\textbf{The two parts.} The \code{entity} is the black-box view: the ports, and nothing about what
happens inside. The \code{architecture} is the implementation behind it. \worth{1}

\textbf{What the split makes possible.}
\begin{itemize}
  \item One design can \textbf{instantiate} another knowing only its entity, without reading the
    architecture.
  \item An implementation can be \textbf{rewritten} without touching anything that depends on it,
    because nothing that depends on it was allowed to see it.
\end{itemize}
\noindent\worth{1}

\textbf{What is different from a C header.} The shape is the same --- a header beside a source
file, or an interface beside its class --- but VHDL \textbf{enforces} it, for every module, with no
way to skip the interface and no way to leak the implementation through it. A C header is a
convention: the compiler does not require one, nothing stops a translation unit from exposing its
internals, and the interface and the implementation can disagree until the linker or the run finds
out. In VHDL the entity is the only way in. \worth{1}

% ==========================================================================================
\answersection{2}{The map, the vector, and the process that lies}{14}

\answerpart{a}{4}
The map, \code{AB} down the rows and \code{CD} across, both in Gray-code order
\code{00, 01, 11, 10}:

\begin{narrowtable}{@{}ccccc@{}}
  \code{AB}\,\textbackslash\,\code{CD} & \code{00} & \code{01} & \code{11} & \code{10} \\
  \midrule
  \code{00} & \textbf{1} & \textbf{1} & & \\
  \code{01} & \textbf{1} & \textbf{1} & \textbf{1} & \textbf{1} \\
  \code{11} & & & \textbf{1} & \textbf{1} \\
  \code{10} & & & \textbf{1} & \textbf{1} \\
\end{narrowtable}
\credit{1 mark for the filled map}

\textbf{A minimal cover, three groups of four:}

\begin{narrowtable}{@{}llll@{}}
  \thead{Group} & \thead{Cells} & \thead{Constant} & \thead{Term} \\
  \midrule
  The top-left block & \code{0000}, \code{0001}, \code{0100}, \code{0101} & \code{A = 0},
    \code{C = 0} & \code{A'C'} \\
  The bottom-right block & \code{1010}, \code{1011}, \code{1110}, \code{1111} & \code{A = 1},
    \code{C = 1} & \code{AC} \\
  The whole \code{AB = 01} row & \code{0100}, \code{0101}, \code{0111}, \code{0110} & \code{A = 0},
    \code{B = 1} & \code{A'B} \\
\end{narrowtable}

\begin{console}
X = A'C' + AC + A'B
\end{console}

{\itshape (3 marks: 1 per group with its cells and term. Nothing here wraps: all three groups are
rectangles you can draw without leaving the map, which is what makes this one easier to see than
the four-corner shapes. \code{A'B} may be replaced by \code{BC} throughout --- see (b).)\par}

\answerpart{b}{2}
\textbf{The colleague's answer.} \code{X = A'C' + AC + A'B + BC} is \textbf{correct}: it computes
the same function, since every group in it is a legal group of cells that are genuinely \code{1}.
It is not minimal.

\textbf{Which group to drop: either \code{A'B} or \code{BC}, but not both.} \code{A'C'} and
\code{AC} are forced --- \code{0000} and \code{0001} sit in no other legal group of four, and
neither do \code{1010} and \code{1011} --- and between them they leave exactly two cells uncovered,
\code{0110} and \code{0111}. Both \code{A'B} and \code{BC} cover that pair, and each of them has its
other two cells already covered by one of the forced groups. So one of the two is redundant and
either one may go.

The rule is step 3 of the procedure: \textbf{use as few groups as will do it, and drop any group
whose cells are every one of them already covered by another.} Nothing in the earlier steps forbids
a redundant group, which is exactly why the step exists --- a redundant group is a redundant
AND-term, correct but not minimal. \worth{1}

\textbf{How many minimal answers.} \textbf{Two.}

\begin{console}
X = A'C' + AC + A'B        and        X = A'C' + AC + BC
\end{console}

\noindent are both three terms of two literals, and both minimal. What that means for marking:
check the \textbf{function}, not the string. Two candidates can hand in different equations and
both be right, and the only way to tell is to evaluate. What is \emph{not} right is keeping both,
which is the colleague's answer. \worth{1}

\answerpart{c}{4}
\textbf{Slicing.} \code{input(7)} drives \textbf{bit 3} of the port. The two are matched \textbf{by
position, left to right}, not by index number: \code{input(7)} drives bit 3, \code{input(6)} drives
bit 2, and so on. The \textbf{lengths} have to match; the \textbf{numbering} does not.
\code{input(7 downto 4)} really does have the index range \code{7 downto 4} and not
\code{3 downto 0}, which is the detail that surprises a software engineer. \worth{2}

\textbf{Why \code{count + 1} does not compile.} \code{std_logic_vector} says nothing about what the
bits \emph{mean}, so it offers only the operations that make sense for a bundle of wires: the
logical operators applied bit by bit, comparison, and concatenation. Arithmetic is not on that list
because nothing has declared whether those four wires are an unsigned number, a signed one, or four
unrelated signals. That is the type doing its job: the meaning genuinely is missing, so the
operation is refused rather than guessed at. \worth{1}

\textbf{What this book uses instead.} A plain integer with a range:

\begin{vhdlcode}
signal count: natural range 0 to 9;
\end{vhdlcode}

\noindent which adds and compares directly, and which synthesis turns into exactly the flip-flops
the range needs. Vectors stay what they are, bundles of wires for things that really are bundles of
wires. The two worlds meet only at a boundary, with the conversions from \code{ieee.numeric_std}:

\begin{vhdlcode}
count_v <= std_logic_vector(to_unsigned(count, 4));   -- integer to vector
count   <= to_integer(unsigned(count_v));             -- vector to integer
\end{vhdlcode}
\noindent\worth{1}

\answerpart{d}{2}
\begin{vhdlcode}
-- One expression.
x <= (a and b) or ((not c) and d);
\end{vhdlcode}

\begin{vhdlcode}
-- The same circuit, with the second term named.
architecture behaviour of combo_logic is
signal cd: std_logic;
begin
    cd <= (not c) and d;
    x  <= (a and b) or cd;
end architecture;
\end{vhdlcode}
\credit{1 mark for both}

\textbf{What synthesis produces: the same circuit, twice.} An internal signal names a wire that
already existed in the first version --- the output of the AND gate computing \code{c'd} --- so
nothing has been added. On the Cyclone V both are one LUT, because a LUT is a lookup table over its
inputs: a function of four or fewer inputs fits in one however it was written, and naming an
intermediate term does not add an input. The answer is the same for both because VHDL describes
hardware, and the two descriptions describe the same hardware.

\textbf{Which to be handed in six months.} At two terms, honestly, either; the named version's
advantage appears as soon as the expression outgrows one line, and the habit is worth having before
you need it. An answer that simply asserts ``the named one is always better'' is worth less than
one that says where the line is. \worth{1}

\answerpart{e}{2}
\textbf{What a simulator does.} The process re-runs only when a signal in its \textbf{sensitivity
list} changes, and the list holds only \code{sel}. So a change on \code{a}, \code{b}, \code{c} or
\code{d} while \code{sel} holds still is invisible: \code{x} keeps whatever value it took the last
time \code{sel} moved. That is not a multiplexer. It behaves like a multiplexer whose output is
captured whenever the selector changes and held in between.

\textbf{What synthesis does.} It builds the multiplexer. The tool reads the logic, works out that
\code{x} depends on \code{sel}, \code{a}, \code{b}, \code{c} and \code{d}, and produces a real
4-to-1 mux --- usually emitting a warning that the sensitivity list is incomplete, and ignoring it
either way. \worth{1}

\textbf{Why they differ.} The sensitivity list is a \textbf{simulation} construct: it tells the
simulator when to re-evaluate the process. It is not hardware, so there is nothing for the synthesis
tool to build from it.

\textbf{The class of bug: the simulator and the chip model different things.} That is the one class
where a green testbench stops meaning anything, because the thing that passed is not the thing that
ships. The book's other example is the opposite direction --- GHDL aborting on a counter wrap the
drawing and the chip both perform happily. The lesson is the same in both: where the two models
disagree, neither one alone is evidence. \worth{1}

Worth a comment rather than a mark: the testbench passes here only because it changes \code{sel}
for every case it checks. One that held \code{sel} still and varied the data would fail
immediately, which is a useful thing to know about your own testbenches.

% ==========================================================================================
\answersection{3}{Storage, and the two ways a value updates}{13}

\answerpart{a}{3}
\textbf{How the two latches are driven.} In opposition. The master's own transparency control is
the inverted clock and the slave's is the clock, so exactly one of the two is ever transparent.

\textbf{Why that gives edge triggering.} While the clock is \textbf{low} the master is transparent
and follows \code{D}, and the slave is closed, holding \code{Q}. At the \textbf{rising edge} the
master closes, freezing whatever \code{D} was at that instant, and the slave opens and passes that
frozen value to \code{Q}. While the clock is \textbf{high} the slave is transparent but its input
cannot change, because the master ahead of it is closed; so \code{D} may do anything it likes and
\code{Q} will not move. At the falling edge the slave closes and the master reopens for the next
capture.

The net effect is that exactly one value passes per period, captured at the instant the clock goes
low to high, and \code{Q} never changes while the clock sits at either level. That external
behaviour is what \code{rising_edge()} models. \worth{2}

\textbf{The latch's weakness.} It is transparent for the \emph{entire} time \code{enable} is high,
so any glitch on \code{D} during that window passes straight through to \code{Q}.

\textbf{Why that is worse in a large system.} In a small circuit you can look at every path feeding
\code{D} and convince yourself it is stable across the phase. In a large one there are thousands,
each with its own delay, and guaranteeing that all of them are settled throughout an entire clock
phase is far harder than guaranteeing they are settled at one instant --- which is all the
flip-flop asks for. \worth{1}

\answerpart{b}{4}
\begin{vhdlcode}
library ieee;
use ieee.std_logic_1164.all;

entity register4 is
    port(clock, reset_n, enable: in  std_logic;
         d                     : in  std_logic_vector(3 downto 0);
         q                     : out std_logic_vector(3 downto 0));
end entity;

architecture behaviour of register4 is
begin
    REGISTER_PROCESS: process(clock, reset_n) is
    begin
        if (reset_n = '0') then
            q <= (others => '0');
        elsif (rising_edge(clock)) then
            if (enable = '1') then
                q <= d;
            end if;
        end if;
    end process;
end architecture;
\end{vhdlcode}

{\itshape (2 marks: 1 for the entity and port order, 1 for the process with \code{clock} and
\code{reset_n} in the sensitivity list, the reset branch outside the clocked branch, and the enable
inside it)\par}

\medskip
\code{q} is only ever assigned here, never read, so no internal signal is needed. A design that
toggles or otherwise reads its own output does need one, because an output port cannot be read
back.

\textbf{Flip-flops: four.} What decides it is the \textbf{width of the signal assigned inside the
clocked process} --- every signal assigned there becomes flip-flops, one per bit. Four bits, one
process, four flip-flops, no different in principle from four \code{d_flip_flop} instances wired in
parallel. \worth{1}

\textbf{Reset branch first: meaning, not style.} The \code{if}/\code{elsif} is read in priority
order, so the reset branch coming first is what makes the reset take priority and act with no clock
edge involved. Put \code{rising_edge(clock)} first and the reset is only examined at an edge, which
is a \textbf{synchronous} reset: a different circuit, one that ignores a reset pulse landing
entirely between two edges. It is also the shape the synthesis tool pattern-matches to infer an
asynchronous reset at all. \worth{1}

\answerpart{c}{4}
\textbf{The trace.} Both are signals, so both assignments \textbf{schedule}: neither takes effect
until the process suspends, and every right-hand side in the pass reads the value the signal held
\emph{before} the edge. \code{stage2_s <= stage1_s;} therefore reads the old \code{stage1_s}, not
the \code{din} the line below is putting there, and the two lines describe two flip-flops in
series.

\begin{narrowtable}{@{}lcccc@{}}
  \thead{Edge} & \thead{\code{din}} & \thead{\code{stage1_s} after} & \thead{\code{stage2_s} after}
    & \thead{\code{q} after} \\
  \midrule
  (after reset) & -- & \code{0} & \code{0} & \code{0} \\
  1 & \code{0} & \code{0} & \code{0} & \code{0} \\
  2 & \code{1} & \code{1} & \code{0} & \code{0} \\
  3 & \code{0} & \code{0} & \code{1} & \code{1} \\
  4 & \code{0} & \code{0} & \code{0} & \code{0} \\
\end{narrowtable}

\noindent The \code{1} on \code{din} at edge 2 reaches \code{stage1_s} at edge 2 and \code{q} at
edge 3: \textbf{two} edges, not one.

{\itshape (2 marks: 1 for the table, 1 for reading each right-hand side against the pre-edge value
rather than asserting the answer)\par}

\medskip
\textbf{The rule.} A signal assignment inside a process \textbf{schedules} a new value, applied only
when the process suspends; every read in the same pass sees the old value. It is the same rule
Chapter~\ref{c3:ch} states for the two-assignment process, and it is what lets a clocked process
describe a chain of flip-flops in whatever order the lines happen to be written.

\textbf{Swapping the two assignments changes nothing.} Write \code{stage1_s <= din;} first and
\code{stage2_s <= stage1_s;} second, and \code{stage2_s} still reads the pre-edge \code{stage1_s},
because scheduling does not care which line came first. That is the point worth taking away:
\textbf{with signals in a clocked process, order does not matter}, which is exactly why the author's
line-order reasoning was the wrong tool. A \code{variable} is the object for which order
\emph{would} matter, and Chapter~\ref{c5:ch} is where the book picks that up. \worth{1}

\textbf{The change that gives one cycle.} Delete \code{stage2_s} and drive the output from the first
stage: \code{stage1_s <= din;} in the process, \code{q <= stage1_s;} outside it. One assignment,
one register.

\textbf{Flip-flops.} The listing as written builds \textbf{two} --- \code{stage1_s} and
\code{stage2_s} are each assigned under \code{rising_edge(clock)} with an asynchronous reset, which
is a flip-flop apiece; \code{q <= stage2_s;} is a wire and costs nothing. After the fix,
\textbf{one}. The extra edge of latency and the extra flip-flop are the same fact seen from two
sides. \worth{1}

\answerpart{d}{2}
\textbf{Why one edge.} Every flip-flop in the design changes at the same well-defined instants, so
the circuit's behaviour is defined at a sequence of points rather than continuously. \worth{1}

\textbf{What that buys.} It turns a continuous-time problem into a discrete one. You can reason
\textbf{one clock cycle at a time} --- ``given the state at this edge, what is the state at the
next?'' --- instead of tracking every signal continuously and worrying about the order in which
things happen in between. Between edges, nothing has to be true except that everything settles
before the next one arrives, which reduces the whole timing question for an arbitrarily large design
to a single number: the longest path against the period. That is what makes a circuit too large to
see all at once analyzable at all, by you or by the tool. \worth{1}

% ==========================================================================================
\answersection{4}{One module, several widths}{13}

\answerpart{a}{5}
\begin{vhdlcode}
library ieee;
use ieee.std_logic_1164.all;

architecture behaviour of button_sync is
-- Three stages, each COUNT bits wide: two synchronizing, one remembering.
signal button_s1_n, button_s2_n, button_s3_n: std_logic_vector(COUNT-1 downto 0);
begin
    -- Falling edge: pressed now, released the cycle before. Bit by bit.
    button_edge_s2 <= (not button_s2_n) and button_s3_n;

    BUTTON_PROCESS: process(clock, reset_s2_n) is
    begin
        if (reset_s2_n = '0') then
            button_s1_n <= (others => '1');
            button_s2_n <= (others => '1');
            button_s3_n <= (others => '1');
        elsif (rising_edge(clock)) then
            button_s1_n <= button_n;
            button_s2_n <= button_s1_n;
            button_s3_n <= button_s2_n;
        end if;
    end process;
end architecture;
\end{vhdlcode}

{\itshape (3 marks: 1 for three vector signals declared in terms of \code{COUNT}, 1 for the
three-stage chain in a clocked process with an asynchronous reset, 1 for the edge expression as a
concurrent assignment outside the process. Deduct nothing for writing the edge inside the process
as a registered pulse, provided the candidate says it arrives a cycle later.)\par}

\medskip
\textbf{Which flip-flops do what.} Stages 1 and 2 are the \textbf{double-flop synchronizer}, so
\code{button_s2_n} is the stable, safe-to-use current value. Stage 3 stores last cycle's
\code{button_s2_n} and is the \textbf{edge detector's own} flip-flop; comparing stage 2 against
stage 3 gives edge detection for free, high for exactly one clock cycle. \worth{1}

\textbf{The reset value: \code{(others => '1')}.} The buttons are active-low, so \code{1} means
released. Reset the chain to \code{0} and it comes up claiming every button is already held down,
which is a lie about the world for the two cycles it takes the real level to propagate --- and,
concretely, a press that is already being held when the reset is released is \textbf{lost}, because
the chain never sees a high-to-low transition on it. This is the same reason Chapter~\ref{c4:ch}'s
CircuitVerse drawing ties \code{preset} to \code{1} on these three flip-flops and leaves it alone on
the reset synchronizer's two. \worth{1}

\textbf{Why \code{reset_s2_n}} --- worth a comment rather than a mark, and the one place a candidate
reveals whether they know why the port is named as it is. This is a \textbf{subcomponent}. A
synchronizer does not synchronize its own reset; \code{reset_sync} does that, and the design that
owns the pin instantiates it. Reading the \code{_s2} as a requirement on the caller is the habit
worth having, because from Chapter~\ref{c6:ch} on most of what you write is a subcomponent.

\answerpart{b}{3}
\textbf{The rule: add a generic when something genuinely varies between instances, not on
principle.} \code{button_sync} is instantiated at \code{COUNT = 1} by \code{led_toggle_sync} and
\code{walking_led}, and at \code{COUNT = 2} by \code{fsm_led}: two widths, one file, and that second
instantiation is what makes the generic worth having. \code{reset_sync} is one bit everywhere.

\textbf{What a \code{WIDTH} on \code{reset_sync} would be.} A parameter nobody ever passes anything
but \code{1} to: pure indirection, and one more thing for a reader to check before they can be sure
what an instance is. \worth{1}

\textbf{Where \code{COUNT} ends up: nowhere.} It is fixed before a single gate is placed, so it does
not appear in the built circuit as logic at all. What it did was decide how many flip-flops to
instantiate and how wide the ports are, and then it was gone --- the same thing that happens to a
\code{for} loop's bounds in Chapter~\ref{c5:ch}. \worth{1}

\textbf{Flip-flop counts.} Three stages, \code{COUNT} bits each: \code{COUNT = 2} gives \textbf{6}
flip-flops (plus two inverters and two AND gates for the edge), \code{COUNT = 1} gives \textbf{3}.
A \code{generic map} itself costs nothing in hardware; it is an elaboration-time decision, not a
runtime one. \worth{1}

\answerpart{c}{2}
\textbf{One mechanism explains both.} The chain samples the button once per clock edge, and whether
that \emph{looks} like debouncing depends entirely on the clock period relative to the bounce time.
\begin{itemize}
  \item In CircuitVerse the period is a full second. Bounce settles in far less than that, so every
    bounce transition happens \textbf{between} two samples and by the next edge the chain reads one
    clean level. The circuit is not debouncing; it is sampling too infrequently to see bounce.
  \item At 50\,MHz the period is 20\,ns, and bounce lasts from under a millisecond to a few tens of
    milliseconds --- many thousands of cycles. The chain faithfully detects \textbf{every} bounce
    transition as its own edge, so one physical press produces several pulses and the LED toggles
    several times.
\end{itemize}
\noindent\worth{1}

\textbf{What a real debounce adds}, either of:
\begin{itemize}
  \item an \textbf{analog RC low-pass filter}, optionally with a Schmitt-trigger buffer, on the
    input before it reaches the pin: the RC pair smooths the fast make/break transitions into one
    gradual edge and the Schmitt trigger's two thresholds turn that back into one clean digital
    transition; or
  \item a \textbf{digital debounce counter}: require the input to hold a new level for a few
    milliseconds' worth of clock cycles before accepting it, which builds the artificially slow
    detection window on-chip without slowing the system clock.
\end{itemize}

\textbf{Why on top rather than instead.} Even a perfectly debounced signal is still asynchronous to
your clock: it can still change inside the setup/hold window and still drive a flip-flop
metastable. Debouncing solves a mechanical problem and synchronizing solves a timing one, and
neither solves the other. \worth{1}

\answerpart{d}{3}
\textbf{\code{din} against \code{reset_n} in \code{seq_detect_mealy}.} \code{din} is
\textbf{synchronous serial data}, already clocked by this clock, so it is not an asynchronous input
and the rule does not apply to it. Synchronizing it anyway would delay it by two clock cycles and
shift its timing relative to everything around it --- in a sequence detector that means the machine
is looking at a different part of the sequence than everything else believes, so the bug is a wrong
answer rather than a slow one. \code{reset_n} comes from a push button and is asynchronous, so it
does get a chain. The rule is about where a signal comes from, not about being cautious. \worth{1}

\textbf{\code{reset_s2_n} against \code{reset_n}.} The port name says which side of the boundary
the module is on:
\begin{itemize}
  \item A \textbf{subcomponent} takes \code{reset_s2_n} and never touches a raw asynchronous input.
    It is a part, and a part is entitled to assume the design around it has done the synchronizing.
    \code{serial_rx8}, \code{timer}, \code{counter}, \code{sipo8}, \code{piso8} and \code{sync} are
    all written this way.
  \item A \textbf{top level} takes \code{reset_n}, because it is the thing holding the pin, and it
    instantiates \code{reset_sync} itself. \code{walking_led}, \code{blinker}, \code{fsm_led} and
    both capstones are written that way.
\end{itemize}
\noindent\worth{1}

\textbf{What goes wrong if a subcomponent is wired straight to a push button.} The reset is then
released asynchronously, so the module's flip-flops leave reset on different clock edges. In
\code{serial_rx8} that is a dozen flip-flops between the shift register and the bit counter, and a
receiver that starts counting a byte one cycle before it starts shifting one is broken by its own
reset --- a fault that appears once in a while, at power-on, and never in a testbench that resets
cleanly. \worth{1}

% ==========================================================================================
\answersection{5}{Two kinds of update, and a speed limit}{13}

\answerpart{a}{3}
\begin{booktable}{@{}p{8.2em}LL@{}}
  & \thead{\code{signal} (\code{<=})} & \thead{\code{variable} (\code{:=})} \\
  \midrule
  Declared & The architecture's declarative part & The process's declarative part \\
  Visible from & Anywhere in the architecture & Only within its own process \\
  Update timing & \textbf{Scheduled}: applied when the process next suspends &
    \textbf{Immediate}: applied before the next statement \\
  Assigned twice in one pass & Only the last survives; earlier ones are discarded & Every one
    takes effect, in order \\
\end{booktable}
\credit{2 marks, or 1 for two of the four rows}

\textbf{The rule of thumb.} Use a \textbf{variable} for a scratch computation consumed entirely
within one pass: loop accumulators, temporary swaps, an intermediate value handed to a signal at the
end. Use a \textbf{signal} for anything another process or the outside world reads, and for
anything that has to hold its value across separate passes --- a counter, a shift register's stored
bits, the previous value in an edge detector. \worth{1}

\answerpart{b}{4}
\begin{vhdlcode}
library ieee;
use ieee.std_logic_1164.all;

entity ones_count8 is
    port(bits: in  std_logic_vector(7 downto 0);
         ones: out natural range 0 to 8);
end entity;

architecture behaviour of ones_count8 is
begin
    COUNT_PROCESS: process(bits) is
        variable count: natural range 0 to 8;
    begin
        count := 0;
        for i in bits'range loop
            if (bits(i) = '1') then
                count := count + 1;
            end if;
        end loop;
        ones <= count;
    end process;
end architecture;
\end{vhdlcode}

{\itshape (3 marks: 1 for a \code{variable} accumulator assigned with \code{:=}, 1 for clearing it
at the top of every activation, 1 for \code{bits} in the sensitivity list and a single assignment
to \code{ones} after the loop)\par}

\medskip
A \code{signal} accumulator here scores none of the three: the assignments schedule rather than take
effect, so only the last of the nine survives and \code{ones} reports the final bit rather than a
count.

\textbf{What the range buys.} Two things, one for each tool. It tells \textbf{synthesis} exactly how
many bits the value needs --- four, for \code{0} to \code{8} --- so the register or the wires are
sized to the value rather than to whatever an unconstrained \code{natural} would give. And it is a
\textbf{promise the simulator checks}: \code{natural range 0 to 8} is range-checked, so GHDL aborts
with a range-check failure at the statement that broke it the instant the count leaves the range.

\textbf{What GHDL would do.} Stop, with an error naming the line, rather than carrying on with a
wrong number. A \code{std_logic_vector(3 downto 0)} would hold \code{9} through \code{15} quite
happily and say nothing, which is the difference between a bug that reports itself and one that
does not. \worth{1}

\answerpart{c}{4}
\textbf{The three delays} the clock period has to cover on a flip-flop / logic / flip-flop path:
\begin{enumerate}
  \item the first flip-flop's \textbf{clock-to-Q} delay, between the edge and its output actually
    changing;
  \item the \textbf{combinational delay} through the LUTs and the wires between the two;
  \item the second flip-flop's \textbf{setup time}.
\end{enumerate}
The \textbf{critical path} is the slowest such path in the design, and \textbf{Fmax} is one over
that total. \worth{2}

\textbf{If the period is too short}, the value has not settled when the second flip-flop samples,
so it captures something still changing --- and the design does the wrong thing \emph{reliably and
mysteriously}, which is worse than doing it intermittently, because it looks like a logic bug.

\textbf{Why adding registers can make a design faster.} Splitting a long combinational expression
so that part of the work happens on one cycle and the rest on the next shortens the longest path,
and the clock period only has to cover the longest path. It trades \textbf{latency} --- more cycles
before an answer comes out --- and some flip-flops, for \textbf{throughput}: a faster clock, and
once the pipeline is full, one result per cycle. That trade is the single most common move in
high-speed digital design. \worth{1}

\textbf{What Quartus needs before it reports an Fmax.} An \code{.sdc} constraints file with a
\code{create_clock} on the 50\,MHz input. Without one the timing analyzer has no period to compare
anything against, so it reports every path as \textbf{unconstrained} and gives no Fmax at all ---
not an optimistic number, no number. Separately, a design that ``fails timing'' is one where a path
turned out slower than the period you did declare, which is a different message and a different
problem.

None of this is part of the book: every design here fits inside 20\,ns with room to spare. Knowing
the step exists and what it needs is what is examinable. \worth{1}

\answerpart{d}{2}
\textbf{An uncleared accumulator keeps growing.} The variable is not a fresh local on each
activation, so a pass that starts by adding to it is adding to the \emph{previous} pass's total.
Over a long run the count climbs without bound: with a range-constrained type it aborts on a range
check as soon as it leaves the range, and with an unconstrained one it simply reports a number that
means nothing. The answer it reports is the running total of every activation since the simulation
began, not a count over \code{bits}. \worth{1}

\textbf{Why a signal for persistent state anyway.} It is possible to build state with a variable,
but its updates are immediate rather than scheduled, so it behaves differently from a signal used
the same way. Question~3(c) is the mirror image: a chain of registers stepping one stage per clock
edge is exactly what scheduled updates give you for free, and a variable there would collapse the
chain. A variable is also invisible outside its own process, so no other logic can read it.
Scheduled updates are what make counters and shift registers work, so if state genuinely has to
persist \emph{and} be read elsewhere, reach for a signal. \worth{1}

% ==========================================================================================
% The title is too long for one line: break it ragged, and before its last phrase.
\answersection{6}{Counting to something that is not a power of two}{12}

\answerpart{a}{4}
\begin{vhdlcode}
library ieee;
use ieee.std_logic_1164.all;

entity counter10 is
    port(clock, reset_s2_n, enable: in  std_logic;
         count                    : out natural range 0 to 9);
end entity;

architecture behaviour of counter10 is
-- Internal count, because an output port cannot be read back.
signal count_s: natural range 0 to 9;
begin
    count <= count_s;

    COUNT_PROCESS: process(clock, reset_s2_n) is
    begin
        if (reset_s2_n = '0') then
            count_s <= 0;
        elsif (rising_edge(clock)) then
            if (enable = '1') then
                if (count_s >= 9) then
                    count_s <= 0;
                else
                    count_s <= count_s + 1;
                end if;
            end if;
        end if;
    end process;
end architecture;
\end{vhdlcode}

{\itshape (3 marks: 1 for the internal signal and the entity, 1 for the clocked process with the
asynchronous reset and the enable, 1 for the explicit comparison and clear)\par}

\medskip
A candidate who writes \code{if (count_s = 9)} rather than \code{>=} loses nothing; \code{>=} is the
defensive form the book's \code{timer} uses.

\textbf{What would be unnecessary at \code{0 to 15}.} The whole comparison:
\code{if (count_s >= 9) then count_s <= 0; else ... end if}, leaving
\code{count_s <= count_s + 1;} on its own.

\textbf{Why exactly those.} \code{15} is \code{2^4 - 1}, so incrementing past it produces a carry
out of the fourth bit with nowhere to go, and the four bits that survive are \code{0000}. The wrap
is a consequence of the register's width --- you can point at the bit it ran out of. \code{9} is
not a power-of-two boundary, so nothing in the hardware returns the count to zero there, and the
comparison has to say so.

Worth noting, and worth a comment rather than a mark: the book writes the comparison out even at
\code{0 to 15}, because a VHDL simulator range-checks \code{natural range 0 to 15} and aborts on the
wrap that the gate network and the chip both perform happily. Anything you intend to verify should
say what it means. \worth{1}

\answerpart{b}{4}
\begin{vhdlcode}
library ieee;
use ieee.std_logic_1164.all;

architecture behaviour of piso8 is
signal shift_reg: std_logic_vector(7 downto 0);
begin
    -- Bit 7 is always the next bit to leave.
    serial_out <= shift_reg(7);

    SHIFT_PROCESS: process(clock, reset_s2_n) is
    begin
        if (reset_s2_n = '0') then
            shift_reg <= (others => '0');
        elsif (rising_edge(clock)) then
            if (load = '1') then
                shift_reg <= parallel_in;                  -- Parallel load.
            elsif (shift = '1') then
                shift_reg <= shift_reg(6 downto 0) & serial_in;  -- Shift toward the MSB.
            end if;
        end if;
    end process;
end architecture;
\end{vhdlcode}

{\itshape (3 marks: 1 for the load branch taking precedence over the shift branch, 1 for the shift
expression in this book's direction, 1 for tapping bit 7 as the serial output)\par}

\medskip
What fills the vacated bit 0 is a wiring decision rather than a property of a PISO. Here it is
\code{serial_in}, so tying that low shifts zeros in behind the byte; \code{walking_led} instead
wires the bit falling off the top back round to it, which is what turns a one-shot pattern register
into an endlessly repeating pattern generator.

\textbf{Which bit leaves first: bit 7, the most significant.} Load \code{"10110010"} and
\code{serial_out} presents \code{1}, \code{0}, \code{1}, \code{1}, \code{0}, \code{0}, \code{1},
\code{0} over the following cycles.

\textbf{The two things that separate it from a SIPO}, the shift expression being identical in both:
\begin{enumerate}
  \item PISO has a \textbf{\code{load} branch}, distinguishing loading a new parallel value from
    shifting the existing one out. SIPO has nothing to load.
  \item PISO \textbf{taps \code{shift_reg(7)}} as a one-bit serial output, where SIPO drives the
    whole vector out in parallel.
\end{enumerate}
``SIPO'' and ``PISO'' name how you wire a shift register up, not two kinds of hardware. \worth{1}

\answerpart{c}{2}
The bit counter now counts \textbf{clock edges} rather than bits, so it loses track of the byte the
moment the stream pauses.

\begin{narrowtable}{@{}cll@{}}
  \thead{Edge} & \thead{What arrives} & \thead{\code{bit_count} after} \\
  \midrule
  1 & bit 7 & 1 \\
  2 & bit 6 & 2 \\
  3 & bit 5 & 3 \\
  4 & (pause) & 4 \\
  5 & (pause) & 5 \\
  6 & (pause) & 6 \\
  7 & (pause) & 7 \\
  8 & (pause) & count was already 7: \textbf{\code{data_ready} pulses}, count clears to 0 \\
\end{narrowtable}

\noindent So \code{data_ready} first pulses on the \textbf{eighth clock edge}, which is the fifth and
last cycle of the pause, five bits before the byte is complete. \worth{1}

\code{data_out} at that moment holds the shift register's contents: the five reset zeros it started
with, followed by the three bits that really did arrive --- \code{"00000"} then bits 7, 6 and 5 of
the byte. And it does not recover: the counter is now permanently out of step with the data, so
every subsequent ``byte'' is a window of eight clock edges rather than eight bits.

That is why \code{shift_enable} has to gate the counter as well as the shift. With both gated, a gap
in the stream \textbf{pauses} the receiver mid-byte rather than corrupting it, and it resumes on
exactly the bit it was waiting for --- which is also what makes the board demonstration possible,
one press per bit. \worth{1}

\answerpart{d}{2}
\begin{booktable}{@{}p{7.4em}L@{}}
  & \thead{What feeds flip-flop \code{i}'s \code{D}} \\
  \midrule
  \textbf{Counter} & The adder's sum bit \code{i}: a function of the counter's own bit \code{i} and
    the carry out of every bit below it. The register's output feeds back into its own input through
    an adder. \\
  \textbf{Shift register} & Flip-flop \code{i-1}'s \code{Q}, straight through, with bit 0's
    \code{D} coming from \code{serial_in}. No arithmetic anywhere. \\
\end{booktable}
\credit{1 mark each}

The difference between the two circuits is therefore not in the flip-flops, which are identical,
and not in the clocking, which is identical. It is entirely in the combinational network between
one stage's \code{Q} and the next stage's \code{D}: an adder in one case, a plain wire in the other.
That is the whole of it, and it is why both are ``technically state machines'' while neither is
what Chapter~\ref{c8:ch} means by one --- the state is a count or a shifted pattern rather than a
label whose meaning you chose.

% ==========================================================================================
\answersection{7}{Turning ticks into time}{10}

\answerpart{a}{3}
\textbf{Cycles.} \code{0.1 s * 50,000,000 cycles/s = 5,000,000} clock cycles.

\textbf{\code{TICK_COUNT}.} The book writes \code{TICK_COUNT = 5_000_000}. The cycle-exact value is
\code{TICK_COUNT = 4_999_999}, because the internal counter runs \code{0} through \code{TICK_COUNT}
inclusive, so a full period is \code{TICK_COUNT + 1} cycles. The round number is 1 cycle long in
5,000,000, which is 0.2\,ppm --- far below the board oscillator's own tolerance, which is why the
book quotes it. \worth{2}

\textbf{Blink frequency: 5\,Hz.} \code{timeout} fires every 100\,ms, so the LED \textbf{toggles} ten
times a second; but one blink is on \emph{and} off, so a full cycle takes two timeouts, 200\,ms,
giving 5\,Hz. The number somebody first writes down is 10\,Hz, which is the timeout rate rather than
the blink rate. \worth{1}

\answerpart{b}{2}
\textbf{A freezing enable means the timer resumes rather than restarts.} Whatever was left of the
period when \code{enable} went low is what is left when it comes back, so the first timeout after
re-entering the state arrives anywhere between one clock cycle and a full period later, depending on
where in its count the timer was frozen. A \textbf{clearing} enable would give a full, predictable
period every time, at the cost of never being able to pause anything.

In \code{fsm_led} the visible consequence is that the first LED toggle after re-entering
\code{STATE_BLINK} can come almost immediately, so the first blink interval is short. \worth{1}

\textbf{Holding \code{enable} high permanently.} Neither better nor worse for correctness, and it
does not fix the short first interval: entering \code{STATE_BLINK} would still catch a free-running
timer at an arbitrary point in its period. What it does buy is a blink phase that is independent of
the state machine entirely; what it costs is a 26-bit counter switching continuously when nothing
is reading it, which is power spent on nothing. The freezing enable is the more useful default
because a design that wants free-running behaviour can always tie \code{enable} high, while a design
that wants to pause cannot un-clear a clearing timer. \worth{1}

\answerpart{c}{3}
\begin{vhdlcode}
library ieee;
use ieee.std_logic_1164.all;

architecture behaviour of blinker is
signal reset_s2_n, timeout: std_logic;
signal enable_always      : std_logic;
signal led_s              : std_logic;
begin
    led           <= led_s;
    enable_always <= '1';

    -- This is a top level: it owns the pin, so it synchronizes the reset itself.
    reset_sync1: entity work.reset_sync
        port map(clock, reset_n, reset_s2_n);

    timer1: entity work.timer
        generic map(TICK_COUNT)
        port map(clock, reset_s2_n, enable_always, timeout);

    -- The timeout pulse is an enable on a flip-flop running on the system clock,
    -- never a clock of its own.
    LED_PROCESS: process(clock, reset_s2_n) is
    begin
        if (reset_s2_n = '0') then
            led_s <= '0';
        elsif (rising_edge(clock)) then
            if (timeout = '1') then
                led_s <= not led_s;
            end if;
        end if;
    end process;
end architecture;
\end{vhdlcode}

{\itshape (3 marks: 1 for \code{reset_sync} instantiated and its \code{reset_s2_n} carried to the
timer, 1 for \code{timer} with the generic passed through and \code{enable} tied high, 1 for the
toggle process using \code{timeout} as an enable rather than as a clock)\par}

\medskip
Tying the literal \code{'1'} directly in the \code{port map} is accepted by GHDL and Quartus and
loses nothing; the named signal is the form that is unambiguously legal in every tool.

\code{led_s} exists because an output port cannot be read back, and this process reads its own
output.

\answerpart{d}{2}
\textbf{Why the synchronizer.} The button is \textbf{asynchronous} to the 50\,MHz clock: a human
decides when it changes, with no relationship to the clock edge, so it can change inside the
setup/hold window of the flip-flop sampling it and drive that flip-flop metastable. That the button
has nothing to do with counting is beside the point --- the rule is about where the signal comes
from, not what it controls. \worth{1}

\textbf{What the design would do without one.} Two things, and the second is the one that matters.
The visible symptom is \textbf{bounce}: the raw contacts make and break several times over the first
millisecond, which at 20\,ns per cycle is many thousands of clock cycles, so the enable would
flicker on and off many times per press and the timer would start, stop and restart. The real fault
is \textbf{metastability}: the enable feeds a 26-bit counter and everything downstream of it, and an
unresolved value reaching several destinations can be read as \code{0} by some and \code{1} by
others in the same cycle, so parts of the design disagree about whether the timer is running.
\worth{1}

% ==========================================================================================
\answersection{8}{The output that arrives a cycle early}{13}

\answerpart{a}{5}
\begin{vhdlcode}
library ieee;
use ieee.std_logic_1164.all;

entity seq_detect_01_mealy is
    port(clock, reset_n: in  std_logic;
         din           : in  std_logic;
         y             : out std_logic);
end entity;

architecture behaviour of seq_detect_01_mealy is
-- STATE_HIGH: the previous bit was '1'. Also the reset state, since the line idles high.
-- STATE_LOW : the previous bit was '0'.
type state_t is (STATE_HIGH, STATE_LOW);

signal reset_s2_n: std_logic;
signal state     : state_t;
begin
    -- Top level: it owns the reset pin. din is synchronous serial data and is NOT
    -- synchronized, exactly as in seq_detect_mealy.
    reset_sync1: entity work.reset_sync
        port map(clock, reset_n, reset_s2_n);

    STATE_PROCESS: process(clock, reset_s2_n) is
    begin
        if (reset_s2_n = '0') then
            state <= STATE_HIGH;
        elsif (rising_edge(clock)) then
            case (state) is
                when STATE_HIGH =>
                    if (din = '0') then
                        state <= STATE_LOW;
                    else
                        state <= STATE_HIGH;
                    end if;
                when STATE_LOW =>
                    if (din = '1') then
                        state <= STATE_HIGH;
                    else
                        state <= STATE_LOW;
                    end if;
                when others =>
                    state <= STATE_HIGH;
            end case;
        end if;
    end process;

    -- Mealy output: the current state AND the current input, combinational,
    -- outside any clocked process.
    y <= '1' when (state = STATE_LOW and din = '1') else '0';
end architecture;
\end{vhdlcode}

{\itshape (3 marks: 1 for a two-value enumerated \code{state_t} with a \code{case} transition
process and an asynchronous reset, 1 for the transitions, 1 for the output as a
\textbf{concurrent} conditional assignment reading both \code{state} and \code{din})\par}

\medskip
\textbf{What the states mean.} \code{STATE_HIGH} is ``the previous bit was \code{1}'',
\code{STATE_LOW} is ``the previous bit was \code{0}''. That is the entire history a rising-edge
detector needs, which is why two states are enough. Either reset state is defensible provided the
candidate says which they chose and why; \code{STATE_HIGH} suits a line that idles high, and means a
first \code{1} after reset is not reported as an edge.

\textbf{The line that makes it Mealy.}

\begin{vhdlcode}
y <= '1' when (state = STATE_LOW and din = '1') else '0';
\end{vhdlcode}

\noindent It reads \code{din}. Move that condition inside the clocked process, or drop \code{din}
from it, and the machine becomes Moore. Note also that it must have an \code{else}: leave one out
and you have described a signal that keeps its old value when no condition holds, which is a latch.
\worth{1}

\textbf{What \code{when others} buys.} In behaviour, nothing: every value of \code{state_t} is
already named above it, so the branch is unreachable in the source and synthesis normally optimizes
it away rather than building recovery logic. What it buys is that the \code{case} stays exhaustive
if a state is added later and a branch is forgotten --- a maintenance property, checked at analysis
time.

What it is \textbf{not} is a safety net against a corrupted state register. Once Quartus re-encodes
the machine, most bit patterns are illegal, and recovering from one needs Quartus's Safe State
Machine setting rather than anything you can write in VHDL. \worth{1}

\answerpart{b}{2}
\textbf{Why one more state.} A Moore output may not read \code{din}, so ``a rising edge just
happened'' has to be recorded in the state itself before the output can report it. That means a
third state, entered on the \code{0} to \code{1} transition, existing purely to hold the output high
for one cycle --- exactly the \code{M_TWO} of the book's own Moore-versus-Mealy comparison in
\secref{c8:a5}.

\textbf{What its output does differently.} It arrives \textbf{one clock cycle later}, and it comes
off a register rather than out of a gate, so it cannot follow a noisy input and, if the output is
registered, cannot glitch. \worth{1}

\textbf{Which to reach for: Moore, by default.} Its output depends on no primary input, so it cannot
follow one; a Mealy output can in principle glitch combinationally if its inputs are not already
clean; and a machine whose output depends on fewer things is a machine with fewer ways to be wrong.
It is what the rest of this book uses. Reach for Mealy only when you specifically need the lower
latency or the smaller state count, and know you are trading a synchronous output for a
combinational one when you do. \worth{1}

Credit an answer that adds that ``Moore'' alone does not mean glitch-free: an output decoded
combinationally from several state bits can still glitch while they skew, and what makes an output
glitch-free is registering it.

\answerpart{c}{4}
\textbf{Cycles per bit and per tick.}

\begin{console}
one bit period  = 50,000,000 / 9600 = 5208.33 clock cycles
one sixteenth   = 5208.33 / 16      =  325.52 clock cycles
\end{console}
\noindent\worth{1}

\textbf{The two candidates.} The timer pulses every \code{TICK_COUNT + 1} cycles, so the target is
\code{TICK_COUNT + 1 = 325.52}, that is \code{TICK_COUNT = 324.52}, which is not a whole number:

\begin{narrowtable}{@{}ccl@{}}
  \thead{\code{TICK_COUNT}} & \thead{Cycles per tick} & \thead{Error against 325.52} \\
  \midrule
  324 & 325 & 0.52 cycles short, \textbf{$-$0.16\%} \\
  \textbf{325} & \textbf{326} & 0.48 cycles long, \textbf{+0.15\%} \\
\end{narrowtable}

\noindent The default is \textbf{325}, and the reason is simply that 326 is the closer of the two:
\code{|326 - 325.52|} is \code{0.48} against \code{|325 - 325.52|} = \code{0.52}. Both would work;
neither is exact, because 9600\,baud does not divide 50\,MHz. \worth{1}

\textbf{Why sample mid-bit.} The transmitter runs off a different physical crystal, so the two
clocks drift apart across a frame, and the receiver's idea of where a bit starts came from one
falling edge it inferred rather than from a clock line it was given. Sampling at the middle leaves
half a bit of margin on \textbf{each} side; sampling near a boundary reads the neighbouring bit as
soon as the two drift at all, or as soon as the line's own edges are anything but instantaneous.
What is being tolerated is frequency error between the two crystals, plus the fraction of a tick of
phase error left over by never stopping the timer. \worth{1}

\textbf{The drift.} Ten bit periods is 160 oversample ticks. The receiver's 160 ticks take
\code{160 * 326 = 52,160} clock cycles; ten true bit periods are \code{10 * 5208.33 = 52,083}
cycles. The difference is about \code{77} cycles, which is \code{77 / 325.52 = 0.24} of an
oversample tick.

The margin it started with is \textbf{8 ticks}, half a bit period. So the sample point has drifted
about \textbf{3\% of its margin} by the stop bit, which is why the answer is comfortable and why
resynchronizing on every start bit is enough: the error never gets ten bit periods to accumulate
over more than once. \worth{1}

\answerpart{d}{2}
\textbf{During the break}, a level test re-arms the instant the previous frame's stop bit is judged.
The line is low, so the receiver reads it as a start bit immediately and marches through
back-to-back all-zero frames for as long as the break lasts. Each of those ends on a low stop bit
and is correctly rejected as a framing error, and \textbf{that part is fine}: nothing is handed
over. \worth{1}

\textbf{At the moment the break ends}, one frame is already in flight, and its stop bit is sampled
after the line has returned high. So it is a perfectly well-formed frame as far as the receiver can
tell: correct framing, high stop bit, and a byte made of nothing handed over as real data.

That last frame is the one that matters because it is the only one that \textbf{passes}. Every
other frame during the break is rejected, so a design that gets this wrong looks correct right up to
the moment a cable is unplugged and plugged back in, and then delivers one plausible byte that was
never sent.

The \textbf{edge} test does not re-arm during the break at all: after the first falling edge there
is not another one until the line has recovered. It costs one flip-flop of history, and that is the
whole difference. \worth{1}

% The marks tag as vhdlbook.sty sets it, for what follows.
